% !TeX program = xelatex
% UTF-8. Standalone: xelatex FUO.tex from this directory, or
% xelatex "pediatrics/ped-infectious diseases/FUO/FUO.tex" from repository root.
% Inclusion: define \FUOInNotebook before \input in a XePersian document.
% Source: Google Drive Notebook/FUO; original JPEG bytes, filename order.
% pic01.jpg = 20261011_120513777.jpg; 1PN9jWtwrd-XLWBqVOSOB9n4xvSUPS10y
% pic02.jpg = 20261011_120522982.jpg; 1pUqwKTDT-79FN2rjFqs6WeLq22cbnUxT
% pic03.jpg = 20261011_120529021.jpg; 1viKdU3YVweCn1eSloU5Gd-vDZshWu2mv
% pic04.jpg = 20261011_120536808.jpg; 1EbIMWOXnTM2D__kbCQuQUo1DlqT95m7y
% pic05.jpg = 20261011_120542415.jpg; 13K888EEvFEhwYnqSR9eu1OZJ07R6yGqh
% pic06.jpg = 20261011_120549154.jpg; 1O-ij51vlCwvkjVH0bYm5vupUsyU8KD83
% pic07.jpg = 20261011_120555485.jpg; 1twl_kGF1Lw4Mn-95VUPB0feay4TasVxi
\ifdefined\FUOInNotebook
  \let\FUOfinish\relax
\else
  \def\FUOfinish{\end{document}}
  \documentclass[12pt,a4paper]{article}
  \usepackage[margin=22mm]{geometry}
  \usepackage{amsmath,amssymb}
  \usepackage{graphicx}
  \usepackage{xepersian}
  \IfFontExistsTF{Amiri}{\settextfont{Amiri}}{\settextfont{DejaVu Sans}}
  \setlatintextfont{DejaVu Serif}
  \setlength{\parindent}{0pt}
  \setlength{\parskip}{5pt}
  \setlength{\emergencystretch}{3em}
  \begin{document}
\fi
\providecommand{\FUOunclear}[1]{\textbf{[خوانش نامطمئن: #1]}}
\providecommand{\FUOmath}[1]{\lr{\(#1\)}}
\providecommand{\FUOpage}[2]{\clearpage\section*{تصویر #1: \lr{\texttt{pic#2.jpg}}}}
\providecommand{\FUOheader}[1]{\textbf{سربرگ و حاشیه:} ۲۲۲ ـ \lr{FUO}؛
  \lr{Subject:} (خالی)؛ \lr{Date: / /}؛ شمارهٔ برگه: #1؛ نشان پایین برگه: \lr{YASHA}.\par}

\begin{center}
{\Large\bfseries \lr{FUO}}\\
\lr{Fever of Unknown Origin}\\
\lr{Nelson 22}\\
۲۲۲
\end{center}
\textbf{یادداشت رونویسی:} متن هفت تصویر، شامل نوشته‌های حاشیه و اختصارات انگلیسی، به ترتیب برگه‌ها رونویسی شده است. فلش‌ها به صورت فهرست یا رابطه بازنویسی شده‌اند. خوانش‌های نامطمئن مشخص شده‌اند. عبارت‌ها و گزاره‌های پزشکی منبع اصلاح نشده‌اند. برای علامت‌های مقایسهٔ دست‌نویس، ترتیب نمادها در موارد مبهم عین تصویر آورده شده است. تصاویر کامل در پیوست قرار دارند.

\FUOpage{۱}{01}
\textbf{سربرگ و حاشیه:} ۲۲۲ ـ \lr{FUO}؛ \lr{Fever of Unknown} و در سطر بعد \lr{Origin}؛
\lr{Subject: Nelson 22}؛ \lr{Date:}؛ شمارهٔ برگه: ۱؛ \lr{YASHA}.

\begin{itemize}
\item \lr{FUO} معمولاً تظاهر ناشایع یک بیماری شایع است و این بیماری‌ها اغلب به سادگی درمان‌پذیرند.
\item \lr{FUO}: تبی که علت آن پس از حداقل ۸ روز بررسی یافت نشده است.
\item \lr{FWLS}: اگر ظرف ۷ روز علتش پیدا نشود، تبدیل به \lr{FUO} می‌شود.
\end{itemize}

\subsection*{اتیولوژی \lr{FUO}}
عفونی، روماتولوژیک، خودالتهابی، آنکولوژیک، نورولوژیک، ژنتیکی، فاکتیشس، ایاتروژنیک.
\FUOunclear{واژهٔ پیش از «ایاتروژنیک» به صورت «فاکتیشس» خوانده می‌شود.}

بیماران بعضی معمولاً علائم دیگری جز تب هم دارند.
\FUOunclear{ترتیب واژه‌ها در این سطر به صورت «بیماران بعضی معمولاً» خوانده می‌شود.}

\subsection*{هر زمان باید احتمال تب دارویی را در نظر بگیریم!}
\begin{itemize}
\item با مصرف همهٔ دارویی! \FUOunclear{حاشیهٔ متصل به عنوان، «با مصرف همهٔ دارویی!» خوانده می‌شود.}
\item معمولاً با علائم دیگری همراه نیست.
\item معمولاً \lr{Sustained Fever} است؛ واژهٔ \lr{Fever} زیر \lr{Sustained} نوشته شده است.
\item قطع دارو منجر به قطع تب ظرف \lr{72 h} می‌شود.
\item تب دارویی به ندرت ۱ ماه طول می‌کشد.
\end{itemize}

در \lr{JIA}، تب به صورت \lr{FUO} تظاهر کلاسیک بیماری است.
\FUOunclear{واژهٔ پس از «تظاهر» شبیه «کلاسیک» است.}

\subsection*{شایع‌ترین علل عفونی که در بیمار \lr{FUO} یافت می‌شوند}
\begin{itemize}
\item انتروکولیت باکتریال (سالمونلوزیس)، توبرکلوز، ریکتزیا، سیفلیس، لایم، خراش گربه.
\item تظاهر آتیپیک و طول‌کشیدهٔ عفونت‌های ویروسی شایع.
\item حاشیهٔ پایین، متصل به مورد قبلی: \lr{EBV}، \lr{CMV}، هپاتیت‌ها، آدنوویروس.
\end{itemize}

\FUOpage{۲}{02}
\FUOheader{۲}
\subsection*{ادامهٔ علل عفونی}
\begin{itemize}
\item هیستوپلاسموز، مالاریا، توکسوپلاسموز.
\item تولارمی، بروسلوز، لپتوسپیروز.
\end{itemize}

خود \lr{AIDS} تب یا \lr{FUO} ایجاد نمی‌کند ولی \lr{FUO} در بیمار \lr{AIDS} شایع‌تر است چون عفونت‌های فرصت‌طلب فرصت رشد دارند.
\FUOunclear{عبارت پیش از «فرصت‌طلب» به صورت «عفونت‌های» خوانده می‌شود.}

\subsection*{بیماری‌های التهابی عامل \lr{FUO}}
\begin{itemize}
\item \lr{JIA} و \lr{SLE}: شایع‌ترین‌ها.
\item کاوازاکی؛ \lr{IBD}.
\end{itemize}

\lr{FUO}، \FUOmath{6<} ماه: گرانولوماتوز، خودالتهابی، خودایمنی.
\FUOunclear{علامت سن در تصویر میان «FUO» و عدد ۶ است؛ ترتیب نوشتهٔ عدد و علامت «۶<» حفظ شده و جهت بالینی آن تفسیر نشده است.}

۹۰\% \lr{FUO}ها علت قابل یافتن دارند که ۵۰\% عفونی، \lr{10--20\%} بیماری‌های کلاژن واسکولار + ۱۰\% آنکولوژیک هستند.
\FUOunclear{عدد بازه در تصویر «10--20\%» خوانده می‌شود.}

\subsection*{تشخیص و ارزیابی بیمار با \lr{FUO}}
\begin{itemize}
\item الگوی تب کمک‌کننده است؛ اگرچه اکثر بیماری‌های عامل \lr{FUO} الگوی مشخصی ندارند.
\item \lr{chill} و شاکینگ‌های مکرر تب مطرح‌کنندهٔ سپتی‌سمی است (فارغ از اینکه علت آن چه باشد).
\end{itemize}
\FUOunclear{عبارت پس از chill شبیه «شاکینگ‌های مکرر تب» است؛ واژهٔ انتهای سطر «سپتی‌سمی» خوانده می‌شود.}

\subsection*{علل \lr{FUO} بر اساس سن}
\FUOmath{4<} سال: معمولاً تنفسی یا ادراری، لوکالیزه (آبسه، استئومیلیت)، \lr{JIA} یا لوکمی.
\FUOunclear{در منبع عدد و علامت سن به صورت «۴<» نوشته شده است.}

\FUOpage{۳}{03}
\FUOheader{۳}
\subsection*{ادامهٔ علل بر اساس سن}
نوجوانان: \lr{IBD}، خودایمنی، لنفوم، توبرکلوز + سالمونلوزیس.

درمان در کودکان با \lr{UTI}: \lr{ABT} خوراکی دارد.

\subsection*{تب در بچه با عدم ایمن‌سازی}
\begin{itemize}
\item \lr{B/C}، \lr{CBC}: بررسی باکتریمی.
\item \lr{ABT} تزریقی به صورت امپریکال اگر:
\begin{itemize}
\item لکوسیتوز: \FUOmath{15{,}000\leq}.
\item \FUOmath{10{,}000\leq ANC}.
\item \FUOmath{0.5<PCT} \lr{ng/ml}.
\end{itemize}
\item تا زمانی که نتیجهٔ کشت آماده شود.
\end{itemize}
\textbf{یادداشت عددها:} در دست‌نویس «۱۵،ـ» و «۱۰،ـ» به صورت اختصاری نوشته شده‌اند؛ ۱۵٬۰۰۰ و ۱۰٬۰۰۰ خوانده شده‌اند. نمادهای مقایسه با همان ترتیب افقی تصویر حفظ شده‌اند.

اندیکاسیون \lr{ABT} اورژانسی: \lr{Sepsis}، مننژیت باکتریال.

\subsection*{تعریف \lr{FUO} و انواع آن}
\begin{itemize}
\item تعریف \lr{FUO}: \FUOmath{38^\circ<T} + \FUOmath{3\,wk<}،
  \FUOmath{2<} ویزیت، \FUOmath{1\,wk<} بستری.
\item \lr{healthcare-associated FUO}: \FUOmath{38^\circ\leq T}،
  \FUOmath{1\,wk<}، حین بستری رخ دهد ولی بدون علل.
\item \lr{Immune deficient FUO}: \FUOmath{38^\circ\leq T}،
  \FUOmath{1\,wk<}، کشت منفی بعد از \lr{48 h}.
\item \lr{HIV-related FUO}: مثل \lr{FUO} کلاسیک، فقط \FUOmath{38^\circ\leq T}.
\end{itemize}
\FUOunclear{انتهای سطر healthcare-associated به صورت «حین بستری رخ دهد ولی بدون علل» خوانده می‌شود. مقایسه‌های زمان عین ترتیب نمادهای تصویر ثبت شده‌اند.}

\subsection*{شرح حال \lr{pica}}
عفونت با توکسوپلاسما گوندی یا توکسوکارا کنیس:
\begin{latin}
Visceral Larva Migrans
\end{latin}

داروهای موضعی هم می‌توانند تب دارویی بدهند: قطرهٔ چشمی.
\begin{latin}
Atropine-Induced Fever
\end{latin}

\FUOpage{۴}{04}
\FUOheader{۴}
\begin{itemize}
\item دیس‌اتونومی خانوادگی (\lr{Riley-Day Syndrome}): هایپرترمی راجعه؛ در بین یهودیان شایع‌تر است.
\item تب مدیترانه‌ای خانوادگی (\lr{FMF})؛ سندرم \lr{hyper-IgD}: هر دو وراثت \lr{AR}.
\end{itemize}
\begin{latin}
TNF receptor-associated periodic syndrome \(\longrightarrow AD\)\\
Muckle-Wells Syndrome \(\longrightarrow AD\)
\end{latin}

\subsection*{\lr{Pseudo-FUO}}
چند عفونت خفیف متوالی که والدین به صورت یک دورهٔ طولانی تب احساس می‌کنند.
\begin{itemize}
\item باید قطعاً \lr{R/O} شود.
\item عفونت اولیه معمولاً کاملاً شناخته می‌شود.
\item تهیهٔ \lr{Fever diary}.
\end{itemize}
\FUOunclear{سطر کوتاه پس از تعریف به صورت «باید قطعاً R/O شود» خوانده می‌شود.}

\subsection*{عدم تعریق علی‌رغم تب شدید}
\begin{itemize}
\item دهیدریشن.
\item انهیدروتیک اکتودرمال دیسپلازی.
\item دیس‌اتونومی خانوادگی.
\item مواجهه با آتروپین.
\end{itemize}

\lr{Red, weeping eyes}: بیماری بافت پیوندی، خصوصاً \lr{PAN}.

\textbf{حاشیهٔ متصل به \lr{PAN}:} پلی‌آرتریت نودوزا.
\FUOunclear{واژهٔ اول در حاشیه شبیه «پلی‌آرتریت» نوشته شده و بدون اصلاح حفظ شده است.}

کنژنکتیویت \lr{palpebral} در بیمار تب‌دار: سرخک، کوکساکی ویروس، توبرکلوز، \lr{IMN}، \lr{Cat Scratch}، لنفوگرانولوم ونروم.
\FUOunclear{نام ویروس شبیه «کوکساکی ویروس» است؛ عبارت آخر در حاشیهٔ پایین «لنفوگرانولوم ونروم» خوانده می‌شود.}

\FUOpage{۵}{05}
\FUOheader{۵}
\begin{itemize}
\item کنژنکتیویت بولبار در کودک با \lr{FUO}: کاوازاکی؛ لپتوسپیروز.
\item خونریزی رتینال ملتحمه: \lr{IE} (اندوکاردیت عفونی).
\item یووئیت در \lr{FUO}: سارکوئیدوز، \lr{JIA}، \lr{SLE}، کاوازاکی، بهجت، واسکولیت.
\item کوریورتینیت در \lr{FUO}: \lr{CMV}، توکسوپلاسموز، سیفلیس.
\end{itemize}
\FUOunclear{عبارت مربوط به خونریزی به صورت «رتینال ملتحمه» خوانده می‌شود؛ واژهٔ پیش از «در FUO» در سطر چهارم «کوریورتینیت» خوانده می‌شود.}

\subsection*{پروپتوز}
تودهٔ اربیت، تیروتوکسیکوز، متاستاز نوروبلاستوم، عفونت اربیت، \lr{GPA} (گرانولوماتوز با پلی‌آنژئیت).

\subsection*{بررسی عروق ناخن}
با استفاده از افتالموسکوپ، ساختار عروق ناخن را بررسی می‌کنیم:
\begin{itemize}
\item \lr{Nail Fold Capillary abn}: درماتومیوزیت جوونایل.
\item اسکلرودرمی سیستمیک.
\end{itemize}
ژل لوبریکانت روی ناخن؛ افتالموسکوپ روی \lr{+4}.

\subsection*{\lr{FUO} در زمینهٔ اختلال هایپوتالاموس}
یک نشانه: ناتوانی عضلات سیلیاری در انقباض به دلیل عدم وجود عضلهٔ منقبض‌کنندهٔ اسفنکتر چشم.
\FUOunclear{نام عضلات در این سطر به صورت «سیلیاری» خوانده می‌شود.}

\subsection*{دیس‌اتونومی خانوادگی}
کاهش اشک، فقدان رفلکس قرنیه، زبان \lr{Smooth} و فقدان \lr{fungiform papillae}.
\FUOunclear{واژهٔ قبل از «اشک» شبیه «کاهش» است.}

\FUOpage{۶}{06}
\FUOheader{۶}
\begin{itemize}
\item کاندیدیاز راجعهٔ دهانی: نشانه‌ای از نقص‌های عملکرد ایمنی، خصوصاً نقص لنفوسیت \lr{T}.
\item \lr{DTR} هایپراکتیو: تیروتوکسیکوز.
\item \lr{Point tenderness} روی یک استخوان: استئومیلیت؛ درگیری مغز استخوان در زمینهٔ متاستاز.
\item تندرنس روی عضلهٔ تراپزیوس: مطرح‌کنندهٔ احتمال آبسهٔ ساب‌دیافراگماتیک.
\item تندرنس جنرالیزهٔ عضلانی: درماتومیوزیت، تریشینوزیس، پلی‌آرتریت، کاوازاکی، مایکوپلاسما.
\end{itemize}
\FUOunclear{واژهٔ فهرست عضلانی پس از تریشینوزیس به صورت «پلی‌آرتریت» خوانده می‌شود.}

\subsection*{معاینه و آزمایش}
\begin{itemize}
\item \lr{TR}: بررسی لنف‌نود پری‌رکتال؛ لنفادنوپاتی یا تندرنس؛ آبسهٔ عمق لگن.
\item \lr{OBT+}: کولیت گرانولوماتوز.
\item \FUOmath{ANC<5{,}000}: بعید است که عفونت باکتریال جدی باشد، جز تب تیفوئید.
\item \FUOmath{10{,}000<PMN}: احتمال یک عفونت باکتریال شدید مطرح است.
\item \FUOmath{30<ESR}: التهاب وجود دارد که لازم است بررسی شود.
\item \lr{ESR} پایین احتمال عفونت یا \lr{JIA} را رد نمی‌کند.
\end{itemize}
\FUOunclear{عبارت TR شبیه «بررسی لنف‌نود پری‌رکتال؛ لنفادنوپاتی یا تندرنس» است. عددهای ANC و PMN در تصویر «5,ـ» و «10,ـ» نوشته شده‌اند؛ به صورت ۵٬۰۰۰ و ۱۰٬۰۰۰ خوانده شده‌اند.}

در \lr{FUO}: نیاز به ۲ مارکر را چک کنیم؛ هر دو با هم ضروری نیستند.
\FUOunclear{سطر آخر به صورت «نیاز به ۲ مارکر را چک کنیم؛ هر دو با هم ضروری نیستند» خوانده می‌شود؛ بخشی از متن روی نشان YASHA افتاده است.}

\FUOpage{۷}{07}
\FUOheader{۷}
\begin{itemize}
\item کشت‌های متعدد برای چند روز در میان: \lr{IE}، استئومیلیت، آبسهٔ عمقی.
\item باکتریمی پلی‌میکروبیال: آلودگی، فاکتیشس، علل \lr{GI}.
\end{itemize}
\FUOunclear{عبارت زمانیِ کشت‌ها شبیه «برای چند روز در میان» است؛ واژهٔ پس از «آلودگی» «فاکتیشس» خوانده می‌شود.}

\subsection*{بررسی \lr{TB}}
\begin{itemize}
\item در کودکان \FUOmath{2<} سال با روش \lr{IGRA}.
\item حاشیهٔ بالای \lr{IGRA}: \lr{INF} \FUOmath{\gamma} \lr{release assay}.
\item بزرگسالان: \lr{TST}.
\end{itemize}
\FUOunclear{عدد و علامت سن به صورت «۲<» حفظ شده است. واژهٔ زیر این سطر به صورت «بزرگسالان» خوانده می‌شود؛ اختصار حاشیه در تصویر INF نوشته شده است.}

بیوپسی در اپروچ تشخیصی \lr{FUO} خیلی جایگاهی ندارد.

\begin{itemize}
\item در کودکان تب و عفونت با هم مترادف نیستند.
\item استثنا: درمان \lr{TB} در بیمار \lr{criticall ill} با شک به \lr{TB} منتشر.
\end{itemize}
\textbf{یادداشت املایی:} عبارت انگلیسی در منبع به صورت \lr{criticall ill} نوشته شده است.

در مرحلهٔ بررسی‌های اولیه تب را پایین می‌آوریم.
\FUOunclear{فعل جملهٔ آخر به صورت «می‌آوریم» خوانده می‌شود؛ «نمی‌آوریم» نیز خوانش محتمل است.}

\clearpage
\section*{پیوست: هفت تصویر کامل منبع}
تمام سربرگ‌ها، حاشیه‌ها، نشانه‌ها و متن دست‌نویس در تصاویر زیر محفوظ‌اند.
% The image directory intentionally preserves the user's "infectioud" spelling.
\IfFileExists{../../ped-infectioud diseases/FUO/images/pic01.jpg}{%
  \def\FUOimagebase{../../ped-infectioud diseases/FUO/images/}}{%
  \def\FUOimagebase{pediatrics/ped-infectioud diseases/FUO/images/}}
\providecommand{\FUOoriginal}[2]{%
  \clearpage\subsection*{اصل تصویر #1: \lr{\texttt{pic#2.jpg}}}%
  \begin{center}
  \includegraphics[width=\linewidth,height=0.87\textheight,keepaspectratio]{\FUOimagebase pic#2.jpg}
  \end{center}}
\FUOoriginal{۱}{01}
\FUOoriginal{۲}{02}
\FUOoriginal{۳}{03}
\FUOoriginal{۴}{04}
\FUOoriginal{۵}{05}
\FUOoriginal{۶}{06}
\FUOoriginal{۷}{07}
\FUOfinish
